% Source: https://artofproblemsolving.com/community/c3h1530675_number_theory_problems
% Stable identity: aops-topic-1530675-q04
A number has exactly $5$ divisors only if it is $p^4$ for a prime $p$ (as $5$ is prime). So $\gcd(m,n)=p^4$, and $p$ appears in $\operatorname{lcm}(m,n)$ with exponent $a\ge4$.

Every prime divisor of $n$ divides $L=\operatorname{lcm}(m,n)$. If $L=p^a q_1^{b_1}\cdots q_k^{b_k}$, then
\[
30=(a+1)(b_1+1)\cdots(b_k+1),\qquad a+1\ge5 .
\]
So $a+1\in\{5,6,10,15,30\}$ and $(b_1+1)\cdots(b_k+1)=\dfrac{30}{a+1}\in\{6,5,3,2,1\}$. A product of $k$ factors each at least $2$ is at least $2^k$, so $2^k\le6$ and $k\le2$; hence $L$, and so $n$, has at most $3$ distinct prime factors.

Three is possible: take $n=2^4\cdot3\cdot5^2=1200$ and $m=2^4=16$. Then $\gcd(m,n)=16$ has $5$ divisors and $\operatorname{lcm}(m,n)=1200$ has $5\cdot2\cdot3=30$ divisors. The answer is $\boxed{3}$.
